% File: quarto/kennsluakademia-template.tex
% Pandoc/Quarto template that maps article.qmd metadata onto kennsluakademia_conf.cls,
% so the Quarto PDF matches the hand-written main.tex.
\documentclass{kennsluakademia_conf}

\conference{Ráðstefna Kennsluakademíu opinberu háskólanna}
\address{Veröld Háskóla Íslands}
\date{20 nóvember, 2026}
\keywords{Lykilorð 1, lykilorð 2, lykilorð 3, lykilorð 4}

\title{Titill}

\author{%
Fyrsti höfundur\autid{1}{0000-0000-0000-0000},
Annar höfundur\autid{1}{0000-0000-0000-0000},
Þriðji höfundur\autid{2}{0000-0000-0000-0000}%
}

\affil{1}{Deild, háskóli}
\affil{2}{Önnur deild, háskóli}

% Helpers pandoc's LaTeX writer expects
\providecommand{\tightlist}{%
  \setlength{\itemsep}{0pt}\setlength{\parskip}{0pt}}
\setlength{\emergencystretch}{3em}
\makeatletter
\newsavebox\pandoc@box
\newcommand*\pandocbounded[1]{%
  \sbox\pandoc@box{#1}%
  \Gscale@div\@tempa{\textheight}{\dimexpr\ht\pandoc@box+\dp\pandoc@box\relax}%
  \Gscale@div\@tempb{\linewidth}{\wd\pandoc@box}%
  \ifdim\@tempb\p@<\@tempa\p@\let\@tempa\@tempb\fi
  \ifdim\@tempa\p@<\p@\scalebox{\@tempa}{\usebox\pandoc@box}%
  \else\usebox{\pandoc@box}%
  \fi%
}
\def\fps@figure{htbp}
\makeatother
\usepackage{longtable,booktabs,array,calc}
\newcounter{none}
\makeatletter
\@ifpackageloaded{caption}{}{\usepackage{caption}}
\AtBeginDocument{%
\ifdefined\contentsname
  \renewcommand*\contentsname{Efnisyfirlit}
\else
  \newcommand\contentsname{Efnisyfirlit}
\fi
\ifdefined\listfigurename
  \renewcommand*\listfigurename{Myndayfirlit}
\else
  \newcommand\listfigurename{Myndayfirlit}
\fi
\ifdefined\listtablename
  \renewcommand*\listtablename{Töfluyfirlit}
\else
  \newcommand\listtablename{Töfluyfirlit}
\fi
\ifdefined\figurename
  \renewcommand*\figurename{Mynd}
\else
  \newcommand\figurename{Mynd}
\fi
\ifdefined\tablename
  \renewcommand*\tablename{Tafla}
\else
  \newcommand\tablename{Tafla}
\fi
}
\@ifpackageloaded{float}{}{\usepackage{float}}
\floatstyle{ruled}
\@ifundefined{c@chapter}{\newfloat{codelisting}{h}{lop}}{\newfloat{codelisting}{h}{lop}[chapter]}
\floatname{codelisting}{Listun}
\newcommand*\listoflistings{\listof{codelisting}{Listanayfirlit}}
\makeatother
\makeatletter
\makeatother
\makeatletter
\@ifpackageloaded{caption}{}{\usepackage{caption}}
\@ifpackageloaded{subcaption}{}{\usepackage{subcaption}}
\makeatother

\begin{document}

\maketitle


\section{Inngangur}\label{inngangur}

Inngangur lýsir vandamáli eða áskorun sem höfundur/ar stóðu frammi fyrir
í kennslu og markmiði verkefnisins.

Kennsluakademían leggur áherslu á að styðja bæði rannsóknarverkefni og
kennsluþróunarverkefni/umbætur á kennslu sem miða að því að efla gæði
náms og kennslu. Akademían samþykkir útdrætti sem byggir á rannsóknum
tengdum háskólum, rannsóknarverkefnum á eigin kennslu (SotL), sem og
útdrætti sem lýsa umbótum á kennsluaðferðum og kennsluþróun sem miða að
því að innleiða nýjar aðferðir, verkfæri eða stefnu til að bæta
námsupplifun nemenda og kennslugæði.

Með þessu vill Kennsluakademían skapa vettvang þar sem fram fer umræða
um hvoru tveggja fræðilegar rannsóknir á kennslu og hagnýtar
kennsluumbætur. Höfundar hvattir til að setja verkefni í fræðilegt
samhengi, með því að vísa í rannsóknir á tengdum sambærilegum verkefnum.

\section{Aðferð}\label{auxf0feruxf0}

Aðferð lýsir framkvæmd verkefnisins, þeirri aðferð eða inngripi sem var
notuð til að leysa vandamálið og hvaða gögnum var safnað til að meta
árangur verkefnisins.

\section{Niðurstöður}\label{niuxf0urstuxf6uxf0ur}

Niðurstöður lýsa árangri af verkefninu, t.d. hvaða áhrif hafði verkefnið
á náms árangur nemenda og upplifun af náminu.

\section{Umræður}\label{umruxe6uxf0ur}

Í umræðum eru niðurstöður ræddar svo sem hvað hefði mátt betur fara og
hvað mögulegu úrbætur væri hægt að gera í framtíðinni. Í umræðum eru
niðurstöður jafnfram ræddar út frá rannsóknium annar á viðfangsefninu.

\nocite{*}
\bibliographystyle{plainnat}
\bibliography{references.bib}


\end{document}
